\section{Related work}
\label{sec:related}

\subsection{Finding out what an RNIC does}

Collie~\cite{collie2022} is the closest relative of the measurement half of this
paper. It searches a high-dimensional configuration space (operation, payload,
maximum transmission unit, queue-pair count, queue depth, batch, gather and
memory-region layout, direction, concurrency, host memory placement) for
configurations in which an RDMA subsystem underperforms, and reports the
resulting anomalies to vendors. We reuse its search dimensions and several of
its published reproducer seeds, and we use its traffic engine, patched. We
differ in three ways. Its anomaly definition assumes a lossless fabric with
priority flow control, so on a fabric without it two of its three detection
channels become invalid near saturation, which our validity review documents
and works around with matched differential controls. Its published cases are
ConnectX-6 and its thresholds are not transferable to a different generation as
constants. And its goal is to find anomalies, whereas ours is to build a model,
so a null result that would end a Collie search (four of the seeds we replayed
produced clean nulls) is for us a positive statement about a block.

Kong et al.~\cite{rnic-micro2023} probe RNIC microarchitectural resources and
their isolation properties, and are the reason our architecture separates
context, translation and work-queue caches rather than treating context
locality as one number. Kalia, Kaminsky and Andersen's design
guidelines~\cite{kalia2016} are the standard reference for the software side of
the interface we model: doorbell batching, inline and BlueFlame publication,
unsignalled completions and the cost of each; FaSST~\cite{fasst2016} and
eRPC~\cite{erpc2019} show how far the datagram path can be pushed and are the
reason our transmit pacer models packets per second separately from bits per
second. Neugebauer et al.~\cite{pcie2018} characterise host PCIe behaviour for
networking and provide the transaction-class accounting our host interface
uses; our own result that PCIe never binds at 100\,GbE is a much weaker claim
than theirs, and holds only at the rates we measured.

SRNIC~\cite{srnic2023} argues for a scalable RNIC architecture with a small
on-chip state footprint, and is the closest published \emph{architecture} to
ours in structure. It differs in intent: SRNIC designs a better NIC, whereas we
reconstruct an existing one, so where SRNIC would remove a limitation we are
obliged to reproduce it.

\subsection{Congestion control}

The device implements DCQCN~\cite{dcqcn2015}, and our reaction point is a
straight implementation of its rate machine. Our contribution to that literature
is negative and specific: on the measured deployment the reaction is 2 to 20
times slower and the recovery about $4.5\times$ slower than the published
defaults suggest, and the notification point's input is not a switch mark at
all. TIMELY~\cite{timely2015} and Swift~\cite{swift2020} replace the marking
signal with delay, and HPCC~\cite{hpcc2019} replaces it with in-network
telemetry; all three are relevant here because the deployment we measured
provides no marking signal, which is precisely the situation those designs
address. IRN~\cite{irn2018} argues that RoCE does not need a lossless network if
the NIC implements better loss recovery, and quantifies the cost of go-back-N
against selective repeat. Our \anom{11} is that cost measured on shipping
silicon: 1.65\,\% packet loss becomes a 26.9\,\% goodput tax at 1\,MiB
messages. Our architecture carries a selective-repeat window as a parameter,
disabled in the ConnectX-5 profile, so that the counterfactual IRN argues for
can be run against a calibrated baseline.

1RMA~\cite{onerma2020} rebuilds the remote-memory interface around
fixed-size, solicited operations with rate control in the host, and is the
clearest statement of the alternative to the connection-oriented model we
reproduce.

\subsection{NIC implementations, open and closed}

Corundum~\cite{corundum2020} is an open-source 100\,Gb/s FPGA NIC and the
reference point for what an open Ethernet datapath on this class of device
costs; it is not an RDMA NIC, so it solves the media and host-interface half of
our problem and none of the transport half. StRoM~\cite{strom2020} is an
open RoCEv2 stack on FPGA with programmable kernels in the data path, and is the
closest existing artifact to the design in \S\ref{sec:arch};
Coyote~\cite{coyote2020} provides the surrounding operating-system abstractions
(memory, networking, scheduling) that a deployment of such a design needs.
NICA~\cite{nica2019} and FlexNIC~\cite{flexnic2016} argue for programmable
inline processing at the network interface, which is orthogonal to our goal but
shares the observation that the interface between host and NIC is the thing
worth designing. Tonic~\cite{tonic2020} is a hardware template for transport
logic at 100\,Gb/s and is the natural implementation vehicle for our B12 and
B13; its segmentation of transport into a small number of reusable modules is
close to our block boundary. The AMD embedded RDMA-enabled NIC~\cite{ernic} is
a commercial RoCEv2 endpoint for the same device family we target, and is the
practical alternative to building this: it is closed, it publishes no timing
contract, and it cannot be calibrated against measured silicon, which is
exactly the gap this work fills.

UCCL~\cite{uccl2025} moves transport into software above the NIC and reports
message-size and outstanding-work anchors on a more recent adapter generation;
we use it as a cross-generation sanity check rather than as a source of
constants. \TODO{confirm the UCCL title, author list and venue or preprint
status before submission.}

\subsection{Where this paper sits}

Everything above either measures an RNIC or builds one. This paper does both
and, more to the point, makes the join explicit: every block in the architecture
names the measurement that constrains it and the evidence class of that
constraint, and every measurement that constrains nothing is either recorded as
a fabric property or retired as a tool artifact. We are not aware of a published
RNIC design that carries its calibration provenance at block granularity, nor of
a published RNIC measurement study whose output is an implementable block
diagram.
